\medskip

Pour chacune des affirmations suivantes, préciser si elle est vraie ou fausse puis justifier la réponse donnée. Toute réponse non argumentée ne sera pas prise en compte.

\bigskip
\Oijk{} est un repère de l'espace.

On considère la droite D qui a pour représentation paramétrique $\left\{\begin{array}{lcl}x&=&\phantom{-}3 - \phantom{3}t \\ y&=&-2 + 3t \\ z&=&\phantom{-}1 + 4t\end{array},\, t \in \R\right.$ et le plan $P$ qui a pour équation cartésienne : $2x -3y + z - 6 = 0$.

\medskip

\begin{enumerate}
\item \textbf{Affirmation :} 
La droite $\mathrm{D}'$, qui a pour représentation paramétrique 

$\left\{\begin{array}{l}x=2+2 t \\ y=4-6 t \\ z=9-8 t\end{array}, t \in \R\right.$, est parallèle à la droite D.

\item On admet que les points A$(-2 ~;~ 3 ~;~ 1)$, B$(1~;~ 3 ~;~-4)$ et C$(6 ~;~ 3 ~;~ 9)$ ne sont pas alignés.

\textbf{Affirmation :} La droite D est orthogonale au plan défini par les trois points A, B et~C.

\item \textbf{Affirmation :} La droite D est sécante avec la droite $\Delta$ qui a pour représentation paramétrique : $\left\{\begin{array}{lcl}x&=&-4+2 t' \\ y&=&\phantom{-}1-3 t' \\ z&=&\phantom{-}2+\phantom{2}t'\end{array} t' \in \R\right.$

\item \textbf{Affirmation :} Le point F$(-3~;~-3~;~3)$ est le projeté orthogonal du point E$(-5~;~0~;~2)$ sur le plan $P$.

\item \textbf{Affirmation :} Il existe exactement une valeur du paramètre réel $a$ telle que le plan $P'$ d'équation $-3x + y - a^{2} z+3=0$ soit parallèle à la droite D.
\end{enumerate}

\smallskip

\textbf{Note du rédacteur} «  pour certaines questions, il est indispensable que le repère soit orthonormé. »

\bigskip
